\documentclass[11pt,letterpaper]{article}
\usepackage[T1]{fontenc}
\usepackage{lmodern,amsmath,amssymb,amsthm,booktabs,microtype}
\usepackage[letterpaper,textwidth=6.25in,textheight=8.52in,top=1.08in,
  headheight=15pt,headsep=18pt,footskip=30pt]{geometry}
\usepackage{fancyhdr}
\usepackage[hidelinks]{hyperref}
\hypersetup{pdftitle={An upper density bound for sets avoiding reciprocal triples},
  pdfauthor={Dmitry Khanukov},pdfsubject={Erdos problem 302}}
\pagestyle{fancy}
\fancyhf{}
\fancyhead[L]{\small\itshape Sets avoiding reciprocal triples}
\fancyhead[R]{\small D. Khanukov}
\fancyfoot[C]{\small\thepage}
\renewcommand{\headrulewidth}{0.2pt}
\newtheorem{theorem}{Theorem}
\newtheorem{lemma}[theorem]{Lemma}
\theoremstyle{definition}
\newtheorem{certificate}[theorem]{Finite certificate}
\newcommand{\N}{\mathbb N}
\newcommand{\cin}{C_{\rm in}}
\newcommand{\cshort}{C_{\rm short}}
\newcommand{\cfull}{C_{\rm full}}
\newcommand{\floor}[1]{\left\lfloor #1\right\rfloor}
\newcommand{\ceil}[1]{\left\lceil #1\right\rceil}
\setlength{\parindent}{1.2em}
\setlength{\emergencystretch}{1.2em}
\raggedbottom
\title{\bfseries An upper density bound for sets\\avoiding reciprocal triples}
\author{Dmitry Khanukov}
\date{Preprint draft with verification supplement \enspace$\cdot$\enspace 20 November 2026}
\begin{document}
\maketitle
\begin{abstract}
Let $f(N)$ be the largest size of a subset of $\{1,\ldots,N\}$ containing
no distinct $a,b,c$ with $1/a=1/b+1/c$. We give a computer-assisted proof of
$\limsup f(N)/N<0.830783022367$. A finite profile bounds the omissions
inside blocks of smooth integers. A packing of disjoint triples on
vertices isolated within these blocks gives further omissions. The
companion Lean development formalizes the counting argument, the
separation of the packings, the checks on the additional dilations,
and the exact arithmetic from the stored ledgers to the stated bound.
The C++ and Python verifiers validate the finite witnesses that
produce the ledgers. A shorter exterior certificate also gives
$\limsup f(N)/N<0.8309$.
\end{abstract}

\section{The extremal problem}
Write $[N]=\{1,\ldots,N\}$. A \emph{reciprocal triple} consists of
three distinct integers $a<b<c$ satisfying $bc=a(b+c)$. Its dilate by
any positive integer is again a reciprocal triple. Erd\H{o}s and
Graham ask how large a set $A\subseteq[N]$ avoiding such triples can
be \cite{EG,Bloom}. Let $f(N)$ denote the maximum.

\begin{theorem}\label{thm:main}
For the extremal function $f$,
\[
  \limsup_{N\to\infty}\frac{f(N)}N
    \ \le\ \cfull\ <\ \frac{830783022367}{10^{12}},\qquad
  \cfull=0.8307830223661779524\ldots .
\]
A shorter finite certificate gives
$\limsup_{N\to\infty}f(N)/N\le\cshort<8309/10000$, with
$\cshort=0.8308999874985765021\ldots$.
\end{theorem}

Cambie's construction gives lower density $5/8$, and van Doorn obtained
the upper bound $9/10$ \cite{Bloom}. The author's earlier two-sided
preprint gives the upper bound $140803024/163562355$ and a lower
bound strictly above $5/8$ for all sufficiently large $N$
\cite{Khanukov}. Kitamura's smooth-block recurrence gives the upper
bound $0.8461739827964010$ \cite{Kitamura}. We use that recurrence
with a new finite profile and count additional triples between blocks.
The proof below uses the version of the finite data that reconstructs
its own initial cover bounds. Determining the extremal density remains open.

\section{Omissions inside smooth blocks}
For a finite $T\subset\N$, a \emph{cover} is a subset of $T$ meeting
every reciprocal triple contained in $T$. Let $\tau(T)$ be the smallest
size of a cover. The complement of a triple-free subset of $T$ is a
cover, so $f(N)=N-\tau([N])$. Deleting all vertices of $T$ that belong
to no triple preserves $\tau(T)$. This observation will allow us to add
the contributions from two families of triples.

\begin{lemma}[Weighted covers]\label{lem:covers}
Let $B_i$ be finite sets with nonnegative integer weights $w_i$. Suppose
$D\ge1$ and $\sum_{i:v\in B_i}w_i\le D$ for every vertex $v\in T$.
Then
\[
  \tau(T)\ \ge\
  \ceil{\frac{1}{D}\sum_i w_i\,\tau(B_i\cap T)}.
\]
\end{lemma}
\begin{proof}
If $H$ covers $T$, then $H\cap B_i$ covers $B_i\cap T$.
After multiplication by $w_i$ and summation, each member of $H$
has contributed at most $D$. Take the minimum over $H$ and use
integrality.
\end{proof}

List the primes through $5003$ in the order $p_1=3$, $p_2=2$,
$p_3=5$, followed by the remaining primes in increasing order.
Put $P_j=\{p_1,\ldots,p_j\}$ and $Q_j=\prod_{i\le j}p_i$, with
$Q_0=1$. Write $S_j(n)$ for the $P_j$-smooth members of $[n]$.
Every $v\in[N]$ factors uniquely as $m s$ with $(m,Q_j)=1$ and $s$
$P_j$-smooth. Thus the sets $mS_j(\floor{N/m})$ partition $[N]$.
At each level there is also a disjoint union
\begin{equation}
 S_j(n)=S_{j-1}(n)\ \sqcup\ p_j S_j(\floor{n/p_j}).
 \label{eq:split}
\end{equation}
Every cover meets the internal triples in both parts. Consequently,
the two corresponding cover lower bounds may be added.

A profile is an integer array $U_j(n)\le\tau(S_j(n))$. Start with
$U_0(n)=U_j(0)=0$. Up to $X=63{,}000{,}000$ take the greatest of
the value at $(j,n-1)$, the inherited certified value, any new finite
cover bound there, and the sum
\[
 U_{j-1}(n)+U_j(\floor{n/p_j}).
\]
Above $X$ use only this last expression. Induction on the level and
then the threshold proves soundness. Define the correction
\begin{equation}
 d_j(n)=U_j(n)-U_{j-1}(n)-U_j(\floor{n/p_j}).
 \label{eq:correction}
\end{equation}
It is nonnegative and vanishes for $n>X$. Store each maximal interval
$[a,b)$ on which $d_j=h>0$ as a cell $(j,a,b,h)$. The corrections
in \eqref{eq:correction} are computed after combining the lower bounds.

\begin{lemma}[Counting a correction]\label{lem:telescoping}
Let $F_j(N)=\sum_{m\le N,\,(m,Q_j)=1} U_j(\floor{N/m})$. Then
\[
 F_j(N)-F_{j-1}(N)
    =\sum_{m\le N,\,(m,Q_{j-1})=1} d_j(\floor{N/m}).
\]
A cell $(j,a,b,h)$ contributes
$h\,\varphi(Q_{j-1})Q_{j-1}^{-1}(1/a-1/b)$ to the limiting
density of these omissions.
\end{lemma}
\begin{proof}
Iterate \eqref{eq:correction} down the powers of $p_j$. Every integer
coprime to $Q_{j-1}$ occurs exactly once as $m p_j^e$ with
$(m,Q_j)=1$. Reindexing the finite sums gives the identity. The cell
counts those $m$ in $(N/b,N/a]$ that are coprime to $Q_{j-1}$.
Each full period of length $Q_{j-1}$ contains $\varphi(Q_{j-1})$
such integers. The contribution divided by $N$ tends to the stated value.
\end{proof}

\begin{certificate}[Internal profile]\label{cert:in}
The full profile gives the density upper bound
\[
 \cin=0.83096432903342704323078172850532877513478108 .
\]
\end{certificate}
The finite computation verifies $89{,}864{,}656$ weighted rows in
four groups, projects them onto $225$ selected prime levels, and
integrates over all $670$ levels. The full profile has $3{,}094{,}631$
positive cells. For $K=10^{22}$, each cell adds
$\floor{K h(b-a)/(ab)}$ to the mass $M_j$ at its level. With $r_0=K$
and $r_j=\floor{r_{j-1}(p_j-1)/p_j}$ the certified value is
\begin{equation}
 \cin=1-K^{-2}\sum_j r_{j-1}M_j. \label{eq:inside}
\end{equation}
Each floor lowers a nonnegative omission count. By
Lemma~\ref{lem:telescoping}, the unrounded cell counts telescope from $F_0=0$ to $F_{670}$.
On each terminal block, the omitted incident vertices alone cover its
internal triples. Summing these disjoint counts gives at least $F_{670}(N)$
omissions, all on incident vertices. Their lower density is therefore at least $1-\cin$. Section~\ref{sec:verification} specifies the local
checks behind the cell table.

\section{An explicit family on isolated vertices}
For $s\ge 2$ let $I(s)$ be the least possible largest entry of a
reciprocal triple containing $s$, and put $I(1)=\infty$. Indeed, two distinct reciprocals with denominators greater than one have sum at most $1/2+1/3$.
For $s\le64$ the values follow by testing $b<c\le64\cdot65$ and
checking whether $bc/(b+c)$ is integral. The triple
$(s,s+1,s(s+1))$ supplies the cutoff for $s\ge2$.

\begin{lemma}[First incidence]\label{lem:incidence}
For every odd prime $\ell$ and $s\ge1$,
\[
 I(\ell s)\ \ge\ \ell\min\{I(s),(\ell+1)/2\}.
\]
\end{lemma}
\begin{proof}
Write a reciprocal triple as a multiple of
$(uv,u(u+v),v(u+v))$, where $u<v$ and $(u,v)=1$.
Indeed, writing $b:c=u:v$ forces $u+v$ to divide the common factor
of $b$ and $c$. If every entry is divisible by $\ell$, division gives
a triple containing $s$. Otherwise $\ell$ divides one of $u,v,u+v$.
In the first two cases $v\ge\ell$, and in the last case $u+v\ge\ell$
and $v\ge(\ell+1)/2$. In all three cases the greatest entry is at
least $\ell(\ell+1)/2$.
\end{proof}

A row consists of three different primes $p,q,r>257$ and three
labels $1\le s_i\le64$. Its vertices have the form $z_i=s_i d_i$, where
$(d_0,d_1,d_2)=(pq,qr,pr)$. Each row is tested for the reciprocal
identity and distinctness. Across rows the pairs $(d_i,s_i)$ are
distinct. Put $P=5003$ and $\lambda=\max_i z_i$. For each $i$, list
the factors of $d_i$ at most $P$ as $\ell_1<\cdots<\ell_m$. Define
\[
 H_i=d_i\min\left(\{I(s_i)\}\cup
     \left\{\frac{\ell_j+1}{2\ell_1\cdots\ell_{j-1}}:
          1\le j\le m\right\}\right),
 \qquad H=\min_i H_i .
\]
The empty product is one. Rows used below have $\lambda<H<\infty$.
Take $B\ge\max(P,p,q,r)$ and let $R_B$ be the product of the primes
at most $B$.

\begin{lemma}[Extra omissions]\label{lem:exterior}
For every $k$ coprime to $R_B$ with $\lambda\le N/k<H$, the triple
$\{kz_i\}$ forces an omission in addition to those counted in
\eqref{eq:inside}. Different admissible row--multiplier pairs use
disjoint vertices.
\end{lemma}
\begin{proof}
Write $d_i=u_i v_i$ with $u_i$ the product of its prime factors at
most $P$. The part of $kz_i$ supported on primes larger than $P$ is
$kv_i$. Iterating Lemma~\ref{lem:incidence}, first removing the largest
prime factor of $u_i$, gives $v_i I(u_i s_i)\ge H_i$. The block of
$kz_i$ therefore has height $\floor{N/(kv_i)}<I(u_i s_i)$, and this
vertex belongs to no internal triple. All three vertices lie in $[N]$,
so at least one isolated vertex must be omitted. From a vertex, its two
least prime factors above $257$, with multiplicity, recover $d_i$. Its
part supported on primes at most $257$ recovers $s_i$. These identify
the row and position, and division recovers $k$.
\end{proof}

For instance, the first stored row uses $(p,q,r)=(4999,5023,5011)$
and $(s_0,s_1,s_2)=(1,2,2)$. Here $2r=p+q$, which gives the
reciprocal identity directly. The largest vertex is $50{,}340{,}506$.
Since $I(2)=6$, the smallest isolation cutoff is $150{,}299{,}934$.
Thus every eligible $k$ in
$(N/150{,}299{,}934,N/50{,}340{,}506]$ contributes one omission.

The density contributed by a row is at least
$\rho(B)(1/\lambda-1/H)$, where
$\rho(B)=\prod_{p'\le B}(1-1/p')$ over primes. For the explicit list,
$B$ is the exact maximum of $P,p,q,r$. The checker rounds $1/\lambda$
down, $1/H$ up and $\rho(B)$ down. For the first $1{,}001{,}000$ rows
the certified gain is exactly
\[
 g_1=\frac{3217076742527054784820507558149578346746501}
                    {5\cdot10^{46}} .
\]
Subtracting $g_1$ from \eqref{eq:inside} gives
$\cshort=0.83089998749857650213508531835416578356784614998$.
The full explicit list has $1{,}400{,}000$ rows and gives the third
bound in Table~\ref{tab:results}.

\section{Additional dilations and the stronger bound}
The larger construction uses $124{,}652{,}035$ base rows with primes
in $(257,160000]$. All prime factors of a label $s_i$ are at most $31$.
Retain $t=1$ for every base row. Additional multipliers $t\le3{,}000{,}000$
have all prime factors in $(31,257]$. Each vertex is now $k t s_i d_i$.
Unique factorization recovers $s_i$, $t$, $d_i$ and $k$ from the
vertex by the same argument as in Lemma~\ref{lem:exterior}. For
counting $k$, the checker may round $B$ up to the end of a group,
which only restricts the permitted multipliers.

For a proposed cap $b$, verify
\begin{equation}
 I(t s)\ge t\min\{I(s),b\} \qquad\text{for every allowed label }s.
 \label{eq:scale}
\end{equation}
This requires only divisors of $(ts)^2$. Indeed, for $x$ on the left
of a reciprocal identity one has $(b'-x)(c'-x)=x^2$ with $b'<c'$.
For $x$ on the right, the other entries are $x-d$ and $x^2/d-x$,
with $d\mid x^2$, $d<x$, and $2d\ne x$. Taking the least height
over these candidates gives $I(x)$. For the left-hand candidate
$x+d,x+x^2/d$ with $d<x$, the right-hand candidate from the same
$d$ has no greater height, except when $d=x/2$. In that case the
left-hand candidate has height $3x$. These observations give the
integer calculation in the supplement.

When all three row primes exceed $P$, use
\[
  \lambda\le N/(kt)<\min\{H,b\min_i d_i\}.
\]
Equation~\eqref{eq:scale} gives isolation. When exactly one row prime
$\ell$ is at most $P$, also impose $2tb\le\ell+1$. If $\ell\mid d_i$,
Lemma~\ref{lem:incidence} gives
$I(\ell t s_i)\ge\ell\min\{I(ts_i),(\ell+1)/2\}$, so the same window
suffices. The calculation uses the lower endpoint of the group containing
$\ell$ in this last test, which is conservative. With two or more row
primes at most $P$, only $t=1$ is used. Each eligible row and scale
contributes the coprime density times the window width in reciprocals,
divided by $t$. All gains use lower rounding.

Write $g_0$ for the gain from all unit dilates and the scaled rows
whose primes all exceed $P$, and $g_m$ for the further gain from
scaling rows with exactly one internal prime. The latter calculation
uses $9{,}581{,}720$ base rows. The exact computation gives
\[
 g_0=0.0001743788090915927106\ldots,\qquad
 g_m=0.0000069278581574980589\ldots.
\]
The resulting coefficient $\cfull=\cin-g_0-g_m$ is exactly
\[
 \cfull=\frac{A\,10^{36}+B_0}{5\cdot10^{72}},\qquad
 \begin{aligned}
 A   &=4153915111830889762306115146791644878,\\
 B_0 &=453404679592858308347342773184262971 .
 \end{aligned}
\]

\begin{table}[ht]
\centering\small
\caption{Upper bounds for $\limsup f(N)/N$, rounded upwards. Each row
starts from the full internal profile.}\label{tab:results}
\begin{tabular}{@{}lr@{}}
\toprule
Certificate & Upper bound \\
\midrule
Internal profile alone & $0.830964329034$\\
First $1{,}001{,}000$ explicit rows & $0.830899987499$\\
All $1{,}400{,}000$ explicit rows & $0.830892297401$\\
Full generated packing with additional scales & $0.830783022367$\\
\bottomrule
\end{tabular}
\end{table}

\paragraph{One period suffices.}
All the cells, rows and scales used for either conclusion are finite.
For an interval $(N/v,N/u]$ counted with modulus $R$, choose $L$
divisible by $Ru$ and $Rv$. For a scaled exterior row the endpoints
in this rule are $t\lambda$ and $tH'$ with $H'$ its truncated upper
endpoint. A positive $L$ common to all intervals exists. At $N=Lj$
every interval is a union of complete periods and its count is exact.
Consequently $f(Lj)\le C Lj$ with $C=\cshort$ or $C=\cfull$.
As $f$ is nondecreasing, every positive $N$ satisfies
\[
 f(N)\le f\bigl(L\ceil{N/L}\bigr)
       \le C L\ceil{N/L}\le C(N+L).
\]
Division by $N$ proves Theorem~\ref{thm:main}. This argument uses
only finiteness of the period.

\section{Scope of the Lean verification}\label{sec:lean}
The Lean declaration \verb|Formal302.upper_0830783022367_of_certificate|
gives the conclusion of Theorem~\ref{thm:main} for the extremal function
in the formulation of Formal Conjectures. Its finite certificate inputs
are the internal correction cells and the exterior windows with their
assigned triples. The hypotheses specify the ranks up to $X$, the
validity and separation of the exterior rows, and the equality of
both computed totals to their recorded numerators. The declaration
has a kernel-checked proof under these hypotheses.

\paragraph{Internal ranks and the limit.}
Lean reconstructs the rank functions recursively from the positive
cells and proves that the rank bounds up to $X$ suffice at every
height. It proves the telescoping identity, the accounting of only
incident vertices, the exact counts on multiples of a common period,
and the passage to the real $\limsup$. The seed replay lemma reduces
local rank bounds to weighted or otherwise certified finite witnesses
and checks of the rank recurrence. The executable \verb|Replay302Cells|
uses the proved rectangle checker to reproduce the masses for all
$3{,}094{,}631$ positive rectangles from the companion data.

\paragraph{Exterior triples.}
Lean proves injectivity of the prime-pair, label and scale encoding
even when the cutoff for the coprime multiplier differs across rows.
It also proves the divisibility criterion for isolation and checks
isolation for the $57$ admissible labels at all $582$ selected
additional scales. Thus the application of the isolated-vertex
argument and the division by each scale have formal proofs.

\paragraph{Exact accounting.}
The kernel checks the arithmetic of the internal mass table, the
$114$ exterior sieve intervals, and the $149$ groups in the mixed
contribution. It also checks the stated exact rational value of
$\cfull$ and its strict comparison with $830783022367/10^{12}$.
The present Lean project leaves the connection from the generated
finite witnesses to the hypotheses of the final declaration open.
In particular, verification of the local cover ranks and of the
complete exterior row stream is performed by the supplied external
verifiers. Section~\ref{sec:verification} describes their checks.

\section{Checking the finite certificates}\label{sec:verification}
The companion archive supplies the tables, the explicit binary list,
the earlier certificate bank, and the C++17 and Python sources needed
to regenerate the witnesses. The following steps explain which
mathematical property is tested at each stage.

\paragraph{Local ranks.}
Each finite support is tested against $bc=a(b+c)$. For a vertex mask
$M$ containing a triple $e$,
$\tau(M)=1+\min_{v\in e}\tau(M\setminus\{v\})$. An edgeless mask has
rank zero. Weighted packings verify the load at every vertex before
applying Lemma~\ref{lem:covers}. Scaling a template preserves its
edges. Intersecting it with each earlier smooth prefix and recomputing
its rank gives a valid bound at that prefix. Branching on inclusion
of a vertex uses the lower of the two conditional demands. Projecting
onto a smaller prefix extends any cover by all deleted vertices and
charges only vertices in the prefix.

\paragraph{Inherited demands.}
The main verifier uses the self-contained bank variant. Four initially
external requests on $71$-smooth prefixes, at thresholds $420,644,806,814$,
have local lower bounds $71,102,125,126$, respectively. Local witnesses
discharge these requests in the reconstruction of the $4{,}106$-node
bank. The full profile also checks $888$ finite cover exclusions by
branching and $23$ finite minima by variable elimination. For the
latter computation, minimize the number of chosen vertices plus the
number of uncovered edges. Any uncovered edge can be met at cost one,
so this minimum equals the cover number. Eliminating a vertex minimizes
over its two possible choices for each assignment of its remaining
neighbours.

\paragraph{Reproduction.}
From the archive root, \texttt{python3 certificate/verify.py} checks the
frozen tables, the directed rounding and the resulting rational numbers.
To verify the actual explicit rows used in the first assertion, run
\verb|python3 certificate/verify.py --shortlist 1001000|.
For full regeneration, run \verb|python3 certificate/verify.py --full work|
on a 64-bit system with Python 3.10 or later and a C++17 compiler
supporting 128-bit integers. This checks the finite witnesses, the
recurrence, both exterior constructions and the final arithmetic.
An independent arithmetic audit and a record of the checks performed
for this draft accompany the supplement. The formalization notes
describe the remaining Lean obligations for these numerical bounds.

\paragraph{Acknowledgment.}
OpenAI tools assisted with the computational investigation, verification
code and preparation of this draft.

\begin{thebibliography}{9}
\bibitem{EG}
P. Erd\H{o}s and R. L. Graham,
\emph{Old and New Problems and Results in Combinatorial Number Theory},
Monographies de L'Enseignement Math\'ematique 28, Universit\'e de
Genève, 1980.

\bibitem{Bloom}
T. F. Bloom, \emph{Erd\H{o}s Problem \#302}, with observations by
S. Cambie and W. van Doorn. \url{https://www.erdosproblems.com/302}.

\bibitem{Khanukov}
D. Khanukov, \emph{Two-sided computer-assisted progress on Erd\H{o}s
Problem 302}, preprint and computational materials, 2026.
\url{https://doi.org/10.5281/zenodo.21966591}.

\bibitem{OEIS}
The OEIS Foundation Inc., \emph{The On-Line Encyclopedia of Integer
Sequences}, entry A390395. \url{https://oeis.org/A390395}.

\bibitem{Kitamura}
K. Kitamura, \emph{A Lean proof of an improved upper bound for Erd\H{o}s
Problem 302}, 2026. \url{https://github.com/KitaKen1/erdos-302-upper-bound}.
\end{thebibliography}
\end{document}
